\chapter{Female Sterilisation}

\section*{Definition and Overview}
Female sterilisation is a permanent method of contraception that blocks or removes the fallopian tubes so that eggs cannot travel to meet sperm. It is also called tubal ligation, tubal occlusion, or, when the tubes are removed, salpingectomy. The procedure is highly effective, with a failure rate of well under 1\%, and it requires no ongoing effort or cost. Sterilisation should be chosen only after careful thought, because reversal is difficult, expensive, and often unsuccessful. It does not protect against sexually transmitted infections, and it does not change a woman's hormones or her menstrual cycle. This chapter explains how the procedure is performed, its effectiveness, and the things to consider before deciding.

\section*{Epidemiology}
Female sterilisation is one of the most commonly used contraceptive methods in the world, used by hundreds of millions of women, though its popularity in Australia has declined in recent decades with the rise of highly effective long-acting reversible contraception such as intrauterine devices and implants. Around 5--10\% of Australian women of reproductive age have been sterilised. A significant proportion of women who choose sterilisation later regret it, with regret more common in younger women, women who are single or who have few children, and those who made the decision under pressure. Salpingectomy, which also reduces ovarian cancer risk, is increasingly performed in preference to simple tubal ligation.

\section*{Aetiology and Pathophysiology}
Sterilisation works by interrupting the physical pathway of the egg.
\begin{itemize}
    \item \textbf{Tubal occlusion:} the fallopian tubes are blocked, clipped, or sealed so the egg released from the ovary cannot meet sperm, and so sperm cannot reach the egg
    \item \textbf{Tubal removal:} salpingectomy removes the fallopian tubes entirely, which also removes the site where most ovarian cancers begin
    \item \textbf{Effect on hormones:} the ovaries are left intact, so hormone production and the menstrual cycle continue normally; sterilisation does not cause menopause
    \item \textbf{Timing of effect:} the method is effective immediately when performed after childbirth (postpartum), but some methods are delayed in effectiveness until the next scan confirms tubal blockage
    \item \textbf{Permanent decision:} reversal requires rejoining the tubes (tubal reanastomosis), which is major surgery with limited success, especially after salpingectomy
    \item \textbf{Alternative:} for men, vasectomy achieves permanent contraception with a simpler, safer procedure
\end{itemize}

\section*{Clinical Features}
The procedure is surgical, and its features relate to preparation and recovery.
\begin{itemize}
    \item Performed under general anaesthesia in most cases, usually as a day procedure or with an overnight stay
    \item Approaches include laparoscopic keyhole surgery through small incisions, a small abdominal cut after childbirth, or hysteroscopic placement of coils without incisions
    \item A quick recovery, with most women back to normal activities within a few days to a week
    \item Mild pain, bruising, and soreness at the incision sites for a few days
    \item Continued normal periods and hormonal cycles after the procedure
    \item No effect on libido or sexual function
    \item The procedure does not protect against sexually transmitted infections
\end{itemize}

\section*{Differential Diagnosis}
The decision to choose sterilisation involves weighing it against other options.
\begin{itemize}
    \item Long-acting reversible contraception: the intrauterine device and implant, which are as effective as sterilisation but reversible
    \item Vasectomy for the male partner, which is simpler, safer, and more easily reversible
    \item Combined and progestogen-only hormonal methods
    \item Barrier methods, including condoms, which also prevent STIs
    \item Fertility awareness methods
    \item The consideration that a woman's circumstances and desires may change, which underlies the need for careful counselling
\end{itemize}

\section*{When to Seek Medical Care}
Women considering sterilisation should have an unhurried consultation with a doctor or family planning clinic to discuss effectiveness, permanence, alternatives, and the small risks of the procedure. This is not an emergency decision.

\textbf{Call 000 (or local emergency services) for:}
\begin{itemize}
    \item Severe abdominal pain, fever, or heavy bleeding after the procedure, which may indicate infection or a complication
    \item Faintness, collapse, or signs of internal bleeding
    \item Severe pain or symptoms that do not settle in the days after surgery
    \item Difficulty breathing or chest pain after anaesthesia
\end{itemize}

\section*{Diagnosis and Investigations}
\begin{itemize}
    \item \textbf{Counselling consultation:} thorough discussion of permanence, effectiveness, failure rate, risks, and alternatives, documented in the consent process
    \item \textbf{Pregnancy test:} to exclude pregnancy before the procedure
    \item \textbf{Pre-operative assessment:} general health review, and blood tests or other checks as needed for anaesthesia
    \item \textbf{Imaging after hysteroscopic methods:} an X-ray or scan at around 3 months to confirm the tubes are blocked before relying on the method
    \item \textbf{No hormonal testing:} sterilisation does not affect hormones, so none is needed
\end{itemize}

\section*{Management}
\begin{itemize}
    \item \textbf{Choice of procedure:} laparoscopic salpingectomy or ligation, postpartum tubal ligation, or hysteroscopic coils, chosen with the surgeon
    \item \textbf{Anaesthesia:} general anaesthesia for most approaches, with day-case or short-stay recovery
    \item \textbf{Pain relief:} simple analgesia for the first few days
    \item \textbf{Wound care:} keeping incisions clean and dry, and reporting any redness, discharge, or fever
    \item \textbf{Resuming activity:} avoiding heavy lifting and strenuous exercise for about a week, and returning to work as advised
    \item \textbf{Use back-up contraception until effectiveness is confirmed}, especially after hysteroscopic methods
    \item \textbf{Continue condoms:} sterilisation does not prevent sexually transmitted infections
\end{itemize}

\section*{Complications}
\begin{itemize}
    \item Risks of anaesthesia and surgery, including bleeding, infection, and damage to nearby organs, which are uncommon
    \item A small failure rate of under 1\%, including a slightly higher risk that any resulting pregnancy is ectopic
    \item Regret, which is more likely in younger women and those who made the decision in difficult circumstances
    \item Chronic pelvic pain in a minority of women after some methods
    \item The difficulty and low success rate of reversal surgery
\end{itemize}

\section*{Prognosis and Outcomes}
Female sterilisation is highly effective and permanent, with a failure rate of less than 1 in 100 women over a lifetime of use, and salpingectomy adds the benefit of reduced ovarian cancer risk. Most women are satisfied with the decision and the method serves them reliably for the rest of their reproductive years. Outcomes are best when the decision is made freely, with full information, and without pressure or regret factors such as young age, relationship instability, or coercion. For women who are certain they have completed their family, it is a safe and dependable choice.

\section*{Prevention and Self-Care}
\begin{itemize}
    \item Take time to decide, and discuss the permanence with a doctor and your partner
    \item Consider whether a reversible long-acting method would meet your needs
    \item Discuss vasectomy as a simpler option if the male partner is in agreement
    \item Have the procedure when you are not under pressure, such as during a difficult relationship or recent birth
    \item Understand the failure rate and the ectopic pregnancy risk
    \item Use condoms after the procedure to prevent sexually transmitted infections
    \item Seek counselling support if you are experiencing regret after the decision
\end{itemize}

\section*{Sources and Further Reading}
The following reputable sources provide further information for healthcare students and patients seeking reliable, current advice about female sterilisation.
\begin{itemize}
    \item Family Planning NSW. \textit{Female sterilisation fact sheet.} https://www.fpnsw.org.au/
    \item Healthdirect Australia. \textit{Female sterilisation.} https://www.healthdirect.gov.au/
    \item The Royal Australian and New Zealand College of Obstetricians and Gynaecologists. \textit{Patient information on tubal occlusion.} https://ranzcog.edu.au/
    \item NHS. \textit{Female sterilisation.} https://www.nhs.uk/conditions/contraception/
    \item Mayo Clinic. \textit{Tubal ligation.} https://www.mayoclinic.org/
    \item World Health Organization. \textit{Family planning.} https://www.who.int/
\end{itemize}
